\section*{Глава 3. Разработка веб-приложения}
\addcontentsline{toc}{section}{Глава 3. Разработка веб-приложения}
\stepcounter{section}

\subsection{Архитектура приложения и структура проекта}

В соответствии с архитектурой, спроектированной в главе 2, приложение
реализовано как клиент-серверная система, состоящая из трёх компонентов:
клиентского приложения на \textbf{React.js}, серверного приложения на
\textbf{Django} с REST API (\textit{Django REST Framework}) и СУБД
\textbf{PostgreSQL}. Для упрощения развёртывания и обеспечения
одинакового окружения на всех этапах разработки компоненты системы
оформлены в виде отдельных контейнеров \textbf{Docker}, объединённых
файлом \texttt{docker-compose.yml} (листинг~\ref{lst:compose}).

\begin{lstlisting}[language=, caption={Конфигурация контейнеров проекта (docker-compose.yml)}, label={lst:compose}]
services:
  db:
    image: postgres:16-alpine
    environment:
      POSTGRES_DB: algo_platform
      POSTGRES_USER: algo_platform
      POSTGRES_PASSWORD: algo_platform

  redis:
    image: redis:7-alpine

  backend:
    build: ./backend
    env_file: ./backend/.env
    depends_on: [db, redis]
    ports: ["8000:8000"]

  frontend:
    build: ./frontend
    ports: ["5173:5173"]
    depends_on: [backend]
\end{lstlisting}

Серверная часть разделена на отдельные Django-приложения, каждое из
которых отвечает за одну группу сущностей логической модели данных,
описанной в главе 2: \texttt{users}, \texttt{problems},
\texttt{submissions}, \texttt{roadmaps}, \texttt{tournaments},
\texttt{messaging} и \texttt{ai\_assistant}. Такое разделение
соответствует принципу единственной ответственности и позволяет
независимо развивать функциональные блоки платформы (задачи, дорожные
карты, турниры, обмен сообщениями и интеллектуальный ассистент).

Взаимодействие между клиентом и сервером построено на двух протоколах:
REST API (CRUD-операции с задачами, дорожными картами, турнирами и
посылками) и WebSocket (модуль \texttt{channels}, используемый для
личных сообщений в реальном времени и описанный в
подразделе~\ref{subsec:frontend}).

\subsection{Реализация модели данных и слоя хранения}

На основе ER-диаграммы (рисунок~\ref{fig:er_diagram}) для каждой таблицы
логической модели данных реализована соответствующая модель Django ORM.
Использование UUID в качестве первичного ключа, выбранное на этапе
проектирования, реализуется полем \texttt{UUIDField} со значением по
умолчанию \texttt{uuid4}. В листинге~\ref{lst:user_model} приведена модель
пользователя, расширяющая стандартную модель \texttt{AbstractUser} полями
\texttt{role}, \texttt{tournament\_rating}, \texttt{streak} и
\texttt{last\_active\_date} из таблицы \texttt{users}.

\begin{lstlisting}[language=Python, caption={Модель пользователя (apps/users/models.py)}, label={lst:user_model}]
class User(AbstractUser):
    class Role(models.TextChoices):
        USER = "user", "Пользователь"
        PREMIUM = "premium", "Премиум-пользователь"
        ADMIN = "admin", "Администратор"

    id = models.UUIDField(primary_key=True, default=uuid.uuid4, editable=False)
    role = models.CharField(max_length=16, choices=Role.choices, default=Role.USER)
    is_blocked = models.BooleanField(default=False)
    tournament_rating = models.IntegerField(default=0)
    streak = models.IntegerField(default=0)
    last_active_date = models.DateField(null=True, blank=True)
    created_at = models.DateTimeField(auto_now_add=True)

    @property
    def is_premium(self) -> bool:
        return self.role in (self.Role.PREMIUM, self.Role.ADMIN)
\end{lstlisting}

Логика обновления серии активных дней (\textit{streak}), описанная в
главе 2, реализована отдельной функцией в модуле \texttt{services.py}
(листинг~\ref{lst:streak}) и вызывается при каждом обращении пользователя
к своему профилю.

\begin{lstlisting}[language=Python, caption={Обновление серии активных дней (apps/users/services.py)}, label={lst:streak}]
def register_activity(user: User, today: date | None = None) -> User:
    today = today or date.today()
    last_active = user.last_active_date

    if last_active == today:
        return user

    if last_active == today - timedelta(days=1):
        user.streak += 1
    else:
        user.streak = 1

    user.last_active_date = today
    user.save(update_fields=["streak", "last_active_date"])
    return user
\end{lstlisting}

Дорожная карта (\texttt{roadmaps}) и её узлы (\texttt{roadmap\_nodes})
образуют дерево за счёт self-referencing внешнего ключа \texttt{parent}
(листинг~\ref{lst:roadmap_node}): у корневых узлов поле \texttt{parent}
равно \texttt{NULL}, у остальных — указывает на родительский узел в той
же таблице.

\begin{lstlisting}[language=Python, caption={Модель узла дорожной карты (apps/roadmaps/models.py)}, label={lst:roadmap_node}]
class RoadmapNode(models.Model):
    id = models.UUIDField(primary_key=True, default=uuid.uuid4, editable=False)
    roadmap = models.ForeignKey(Roadmap, on_delete=models.CASCADE, related_name="nodes")
    parent = models.ForeignKey(
        "self", on_delete=models.CASCADE, related_name="children",
        null=True, blank=True,
    )
    title = models.CharField(max_length=255)
    description = models.TextField(null=True, blank=True)
    position = models.PositiveIntegerField(default=0)
    created_at = models.DateTimeField(auto_now_add=True)
\end{lstlisting}

Для остальных таблиц (\texttt{problems}, \texttt{submissions},
\texttt{tournaments}, \texttt{tournament\_participants},
\texttt{messages}) реализованы аналогичные модели с полями, типами и
ограничениями, перечисленными в таблицах~2.1--2.7 главы~2. На основе
моделей сгенерированы и применены миграции базы данных PostgreSQL
(\texttt{python manage.py makemigrations} / \texttt{migrate}), что
завершает выполнение задачи 2.3 «Проектирование модели данных».

\subsection{Реализация системы авторизации и ролевой модели доступа}

Для аутентификации пользователей используется протокол \textbf{JWT}
(\textit{JSON Web Token}), реализованный библиотекой
\texttt{djangorestframework-simplejwt}. При успешном входе сервер выдаёт
пару токенов (\texttt{access} и \texttt{refresh}); access-токен передаётся
в заголовке \texttt{Authorization} каждого запроса к API, а при его
истечении клиент автоматически обновляет токен через
\texttt{/api/auth/token/refresh/} (листинг~\ref{lst:api_client}).

\begin{lstlisting}[language=Python, caption={Перехват 401-ответа и обновление access-токена (src/api/client.js)}, label={lst:api_client}]
client.interceptors.response.use(
  (response) => response,
  async (error) => {
    const { config, response } = error;
    if (response?.status === 401 && !config._retried) {
      const refresh = localStorage.getItem("refresh_token");
      if (refresh) {
        const { data } = await axios.post("/api/auth/token/refresh/", { refresh });
        localStorage.setItem("access_token", data.access);
        config._retried = true;
        config.headers.Authorization = `Bearer ${data.access}`;
        return client(config);
      }
    }
    return Promise.reject(error);
  }
);
\end{lstlisting}

Ролевая модель доступа (\texttt{user}, \texttt{premium}, \texttt{admin}),
описанная в подразделе 2.1, реализована на уровне поля \texttt{role}
модели пользователя и набора классов разрешений DRF
(листинг~\ref{lst:permissions}). Класс \texttt{IsAuthorOrReadOnly}
разрешает изменение и удаление задач и дорожных карт только их автору
либо администратору, а класс \texttt{IsPremium} ограничивает доступ к
ИИ-ассистенту премиум-пользователями и администраторами.

\begin{lstlisting}[language=Python, caption={Классы разрешений ролевой модели (apps/users/permissions.py)}, label={lst:permissions}]
class IsPremium(BasePermission):
    """Доступ только для премиум-пользователей и администраторов."""

    def has_permission(self, request, view) -> bool:
        return bool(
            request.user
            and request.user.is_authenticated
            and request.user.is_premium
        )


class IsAuthorOrReadOnly(BasePermission):
    """Изменять и удалять объект может только его автор или администратор."""

    def has_object_permission(self, request, view, obj) -> bool:
        if request.method in SAFE_METHODS:
            return True
        if request.user.role == request.user.Role.ADMIN:
            return True
        return getattr(obj, "author_id", None) == request.user.id
\end{lstlisting}

\subsection{Реализация модуля проверки решений и ИИ-ассистента}

Отправка решения пользователем обрабатывается эндпоинтом
\texttt{/api/submissions/}: созданная запись \texttt{Submission}
передаётся в модуль проверки решений (\textit{judge}), который
сопоставляет результат выполнения программы с закрытым набором тестов
задачи и заполняет поля \texttt{tests\_passed}, \texttt{verdict} и
\texttt{is\_accepted}. Если посылка отправлена в рамках турнира, для
неё дополнительно пересчитывается турнирный балл по формуле~2.1 и
итоговый рейтинг участника по формуле~2.2 (листинг~\ref{lst:scoring}).

\begin{lstlisting}[language=Python, caption={Расчёт турнирного балла и рейтинга (apps/tournaments/services.py)}, label={lst:scoring}]
def submission_score(submission: Submission) -> float:
    """S = T_passed / T_total * 100"""
    if submission.problem.total_tests == 0:
        return 0.0
    return submission.tests_passed / submission.problem.total_tests * 100


def recalculate_rating(participant: TournamentParticipant) -> TournamentParticipant:
    """R = sum(S_j) -- лучший балл по каждой задаче"""
    submissions = Submission.objects.filter(
        tournament=participant.tournament, user=participant.user
    )
    best_score_per_problem = {}
    for submission in submissions:
        score = submission_score(submission)
        best_score_per_problem[submission.problem_id] = max(
            best_score_per_problem.get(submission.problem_id, 0), score
        )

    participant.total_score = round(sum(best_score_per_problem.values()))
    participant.save(update_fields=["total_score"])
    return participant
\end{lstlisting}

Интеллектуальный анализ решений (задача 1.6 «Реализовать интеллектуальный
анализ решений») реализован в виде интеграции с \textbf{DeepSeek API}.
Сервис \texttt{ai\_assistant.services.ask\_assistant} формирует запрос,
включающий условие задачи, текущий код пользователя (если он указан) и
системную инструкцию, запрещающую ассистенту раскрывать готовое решение
(листинг~\ref{lst:ai_service}). Доступ к эндпоинту \texttt{/api/ai-assistant/ask/}
ограничен классом разрешения \texttt{IsPremium}.

\begin{lstlisting}[language=Python, caption={Запрос к ИИ-ассистенту (apps/ai\_assistant/services.py)}, label={lst:ai_service}]
SYSTEM_PROMPT = (
    "Ты -- ассистент по подготовке к олимпиадам по программированию. "
    "Направляй пользователя к решению задачи наводящими вопросами и "
    "подсказками, не приводя готовый код и не раскрывая решение напрямую."
)

def ask_assistant(problem_description: str, user_code: str | None, question: str) -> str:
    messages = [{"role": "system", "content": SYSTEM_PROMPT}]
    context = f"Условие задачи:\n{problem_description}"
    if user_code:
        context += f"\n\nТекущий код пользователя:\n{user_code}"
    messages.append({"role": "user", "content": context})
    messages.append({"role": "user", "content": question})

    response = requests.post(
        settings.DEEPSEEK_API_URL,
        headers={"Authorization": f"Bearer {settings.DEEPSEEK_API_KEY}"},
        json={"model": "deepseek-chat", "messages": messages},
        timeout=30,
    )
    response.raise_for_status()
    return response.json()["choices"][0]["message"]["content"]
\end{lstlisting}

\subsection{Разработка клиентской части}
\label{subsec:frontend}

Клиентская часть реализована на \textbf{React.js} с использованием
\textit{Vite} и маршрутизации \texttt{react-router-dom}. Структура
страниц соответствует функциональным блокам, описанным в подразделе 2.2:
список и страница задачи со встроенным редактором кода, дорожные карты,
турниры и личные сообщения.

Страница задачи объединяет три компонента: условие задачи, встроенную
IDE на базе \texttt{@monaco-editor/react} с выбором языка
(Python, C++, C, Go, Java) и чат ИИ-ассистента. Компонент
\texttt{AIAssistantChat} (листинг~\ref{lst:ai_chat}) скрывает форму
запроса для пользователей с ролью \texttt{user}, реализуя на клиенте то
же ограничение доступа, что и серверный класс \texttt{IsPremium}.

\begin{lstlisting}[language=JavaScript, caption={Компонент чата с ИИ-ассистентом (src/components/AIAssistantChat.jsx)}, label={lst:ai_chat}]
const AIAssistantChat = ({ problemId, code }) => {
  const { user } = useAuth();
  const [question, setQuestion] = useState("");
  const [history, setHistory] = useState([]);

  if (!user || user.role === "user") {
    return (
      <div className="ai-assistant ai-assistant--locked">
        ИИ-ассистент доступен только премиум-пользователям.
      </div>
    );
  }

  const sendQuestion = async () => {
    setHistory((prev) => [...prev, { role: "user", content: question }]);
    const { data } = await client.post("/ai-assistant/ask/", {
      problem: problemId,
      question,
      code,
    });
    setHistory((prev) => [...prev, { role: "assistant", content: data.answer }]);
    setQuestion("");
  };
  // ...
};
\end{lstlisting}

Дорожная карта на клиенте отображается рекурсивным компонентом
\texttt{RoadmapTree} (листинг~\ref{lst:roadmap_tree}), который повторяет
структуру дерева узлов \texttt{roadmap\_nodes}: каждый узел может
содержать список дочерних узлов \texttt{children}, полученных от сервера
в виде вложенного JSON.

\begin{lstlisting}[language=JavaScript, caption={Рекурсивный рендеринг дерева дорожной карты (src/components/RoadmapTree.jsx)}, label={lst:roadmap_tree}]
const RoadmapNode = ({ node }) => (
  <li className="roadmap-tree__node">
    <span>{node.title}</span>
    {node.children?.length > 0 && (
      <ul>
        {node.children.map((child) => (
          <RoadmapNode key={child.id} node={child} />
        ))}
      </ul>
    )}
  </li>
);
\end{lstlisting}

Личные сообщения реализованы поверх WebSocket-соединения: при открытии
диалога клиент сначала запрашивает историю сообщений через REST API
(\texttt{/api/messages/?with=<id>}), а затем подключается к каналу
\texttt{/ws/chat/<id>/}. На сервере соединение обрабатывается
асинхронным консьюмером \texttt{ChatConsumer} на базе \texttt{Django
Channels}: оба собеседника подключаются к одной группе, идентификатор
которой формируется из отсортированной пары их \texttt{id} (листинг~\ref{lst:chat_consumer}), что
обеспечивает доставку сообщений в реальном времени без дополнительной
адресации.

\begin{lstlisting}[language=Python, caption={WebSocket-консьюмер личных сообщений (apps/messaging/consumers.py)}, label={lst:chat_consumer}]
class ChatConsumer(AsyncWebsocketConsumer):
    async def connect(self):
        user = self.scope["user"]
        if not user.is_authenticated:
            await self.close()
            return

        self.other_user_id = self.scope["url_route"]["kwargs"]["user_id"]
        self.room_group_name = "chat_" + "_".join(
            sorted([str(user.id), str(self.other_user_id)])
        )
        await self.channel_layer.group_add(self.room_group_name, self.channel_name)
        await self.accept()

    async def receive(self, text_data):
        data = json.loads(text_data)
        message = await self._save_message(
            self.scope["user"].id, self.other_user_id, data["content"]
        )
        await self.channel_layer.group_send(
            self.room_group_name,
            {"type": "chat_message", "sender_id": str(message.sender_id),
             "content": message.content, "sent_at": message.sent_at.isoformat()},
        )
\end{lstlisting}

Таким образом, на текущем этапе реализованы базовая структура клиентского
и серверного приложений, модель данных и слой хранения в PostgreSQL,
ролевая модель доступа на основе JWT, модуль обработки посылок с расчётом
турнирного рейтинга и интеграция ИИ-ассистента, что соответствует задачам
2--5, сформулированным в главе 1.
